% =============================================================================
% template.tex  —  Cryptography / Mathematics Research Paper
%
% COMPILATION (recommended):
%   pdflatex template  &&  biber template  &&  pdflatex template  &&  pdflatex template
% Or with latexmk:
%   latexmk -pdf -pvc template
%
% Install crypto-math.sty once to your local TEXMF tree:
%   mkdir -p ~/texmf/tex/latex/local
%   cp crypto-math.sty ~/texmf/tex/latex/local/
%   mktexlsr                      # or: texhash
% =============================================================================

\documentclass[11pt, a4paper]{article}

% --------------------------------------------------------------------------- %
%  Load the master style sheet.
%  Pass options:
%    draft     -- disables microtype expansion, marks hyperref draft
%    anonymous -- suppresses author info (double-blind submissions)
% --------------------------------------------------------------------------- %
\usepackage{crypto-math}
% \usepackage[draft]{crypto-math}
% \usepackage[anonymous]{crypto-math}

% --------------------------------------------------------------------------- %
%  Bibliography file(s) — uncomment and adjust path.
% --------------------------------------------------------------------------- %
% \addbibresource{refs.bib}

% --------------------------------------------------------------------------- %
%  PDF metadata  (hyperref)
% --------------------------------------------------------------------------- %
\hypersetup{
  pdftitle  = {Paper Title},
  pdfauthor = {Tess Dore},
  pdfsubject= {Cryptography},
  pdfkeywords={keyword1, keyword2, keyword3},
}

% ============================================================================
%  TITLE BLOCK
% ============================================================================
\title{%
  \textbf{Paper Title}\\[0.5ex]
  \large Subtitle (if any)
}

\author{%
  Tess Dore\\
  \small RippleX Research\\
  \small \texttt{tdore@ripple.com}
}

\date{\today}   % or \date{} to suppress

% ============================================================================
%  DOCUMENT
% ============================================================================
\begin{document}

\maketitle
\thispagestyle{plain}   % page number only on title page

% --------------------------------------------------------------------------- %
\begin{abstract}
  Write a concise abstract here.
  State the problem, your main result, and the key technique.
\end{abstract}

\paragraph{Keywords.} Keyword one; keyword two; keyword three.

% --------------------------------------------------------------------------- %
\tableofcontents
\newpage

% ============================================================================
%  §1  INTRODUCTION
% ============================================================================
\section{Introduction}
\label{sec:intro}

A forward reference works: see \cref{sec:prelim} for background.
Cite a paper like this~\cite{example2024}.

% ============================================================================
%  §2  PRELIMINARIES
% ============================================================================
\section{Preliminaries}
\label{sec:prelim}

\subsection{Notation}

Let $\secpar \in \N$ be the security parameter.
We write $x \getsu S$ for $x$ sampled uniformly from a set $S$,
and $x \getsu \D$ for $x$ sampled from distribution $\D$.
All algorithms are implicitly given $\secp$ as input.
$\negl(\secpar)$ denotes a negligible function.
We use $\bits^n$ for the set $\{0,1\}^n$.

\subsection{Definitions}

\begin{definition}[Hard-core predicate]
  \label{def:hardcore}
  A function $b \colon \bits^* \to \bits$ is a \emph{hard-core predicate} for a
  one-way function $f$ if \ldots
\end{definition}

\begin{definition}[$(\varepsilon, \delta)$-Differential Privacy]
  A randomised algorithm $\mathcal{M}$ is $(\varepsilon,\delta)$-differentially
  private if for all neighbouring datasets $D, D'$ and all $S \subseteq \mathrm{Range}(\mathcal{M})$:
  \[
    \Pr{\mathcal{M}(D) \in S} \;\le\; e^{\varepsilon}\,\Pr{\mathcal{M}(D') \in S} + \delta.
  \]
\end{definition}

% ============================================================================
%  §3  MAIN RESULTS
% ============================================================================
\section{Main Results}
\label{sec:results}

\begin{theorem}[Main theorem]
  \label{thm:main}
  Assuming the $\LWE_{n,q,\chi}$ problem is hard,
  the scheme $\Pi$ in \cref{con:main} is $\mathrm{IND\text{-}CPA}$-secure.
\end{theorem}

\begin{construction}
  \label{con:main}
  The scheme $\Pi = (\KeyGen, \Enc, \Dec)$ is defined as follows.
  \begin{description}[leftmargin=2em]
    \item[$\KeyGen(1^\secpar)$:] Sample $\mathbf{A} \getsu \Zq^{m \times n}$, \ldots
    \item[$\Enc(\pk, \mu)$:] \ldots
    \item[$\Dec(\sk, \ct)$:] \ldots
  \end{description}
\end{construction}

\begin{proof}[Proof of \cref{thm:main}]
  We reduce to $\LWE$.
  Let $\A$ be a $\PPT$ adversary against $\Pi$.
  We construct a reduction $\B$ against $\LWE$ \ldots
  \[
    \Adv{\mathrm{IND\text{-}CPA}}{\Pi}(\A) \;\le\;
    2 \cdot \Adv{\LWE_{n,q,\chi}}{\B}(\A) + \negl(\secpar).
  \]
\end{proof}

\begin{lemma}
  \label{lem:key}
  For all $\mathbf{s} \in \Zq^n$ and $\chi$-distributed $e$,
  \[
    \norm{\mathbf{A}^\top\mathbf{s} + \mathbf{e} \bmod q}_\infty
    \;\le\; q/4
  \]
  with probability $1 - \negl(\secpar)$.
\end{lemma}

% ============================================================================
%  §4  ALGORITHMS
% ============================================================================
\section{Algorithm Example}
\label{sec:alg}

\begin{algorithm}[H]
  \caption{Example algorithm}
  \label{alg:example}
  \begin{algorithmic}[1]
    \Require{Public matrix $\mathbf{A} \in \Zq^{m \times n}$, query $\mathbf{b} \in \Zq^m$}
    \Ensure{Bit $b \in \bits$}
    \State $\mathbf{c} \gets \mathbf{A}\mathbf{A}^\top\mathbf{b} \bmod q$
    \For{$i \gets 1$ to $m$}
      \If{$c_i > q/2$}
        \State $c_i \gets c_i - q$
      \EndIf
    \EndFor
    \State \Return $\lfloor 2\norm{\mathbf{c}} / q \rceil \bmod 2$
  \end{algorithmic}
\end{algorithm}

% ============================================================================
%  §5  CONCLUSION
% ============================================================================
\section{Conclusion}
\label{sec:conclusion}

\ldots

% ============================================================================
%  BIBLIOGRAPHY
% ============================================================================
% With biblatex:
% \printbibliography

% With classic BibTeX (use instead of biblatex block in sty):
% \bibliographystyle{alpha}
% \bibliography{refs}

% ============================================================================
%  APPENDIX
% ============================================================================
\begin{appendices}

\section{Proof of \texorpdfstring{\cref{lem:key}}{Lemma~\ref{lem:key}}}
\label{app:proof-lem-key}

\begin{proof}
  \ldots
\end{proof}

\end{appendices}

\end{document}
